%! AP Statistics — Unit 4, Lesson 4.11 — Group Activity
\documentclass[10pt]{article}
\usepackage{apstats-article}
\usepackage{apstats-boxes}

\begin{document}

\pageheader{Unit 4, Lesson 4.11}{Group Activity: Binomial Parameters in Context}
\namepartnerperiod

\vspace{0.1in}

\begin{scenariobox}[Directions]{sky}
\small
Work through all three tiers. Show all formula setup before computing. Write
complete-sentence interpretations where asked --- the AP exam requires this.
\end{scenariobox}

\vspace{0.1in}

% ── TIER R ────────────────────────────────────────────────────────────────────
\begin{tcolorbox}[colback=goldbg, colframe=goldacc, boxrule=1pt, arc=2mm,
    title={\bfseries Tier R \quad Remediate: Computing Mean and SD},
    fonttitle=\small]
\small

For each binomial distribution, compute $\mu_X$ and $\sigma_X$, then write a
one-sentence interpretation of $\mu_X$.

\vspace{10pt}
\textbf{(a)} $X \sim B(25, 0.40)$

\vspace{4pt}
$\mu_X = np = $ \blank{5cm} $= $ \blank{2cm}

$\sigma_X = \sqrt{np(1-p)} = \sqrt{\blank{3cm}} \approx$ \blank{2cm}

Interpretation of $\mu_X$: \writeline\\[1pt]\writeline

\vspace{12pt}
\textbf{(b)} $X \sim B(100, 0.08)$

\vspace{4pt}
$\mu_X = $ \blank{5cm} $= $ \blank{2cm}

$\sigma_X = \sqrt{\blank{3cm}} \approx$ \blank{2cm}

Interpretation of $\mu_X$: \writeline\\[1pt]\writeline
\vspace{4pt}
\end{tcolorbox}

\vspace{0.1in}

% ── TIER A ────────────────────────────────────────────────────────────────────
\begin{tcolorbox}[colback=sky, colframe=navy, boxrule=1pt, arc=2mm,
    title={\bfseries Tier A \quad Approaching Proficiency: A Full Binomial Analysis},
    fonttitle=\small]
\small

\textbf{Context:} A call center reports that 22\% of customer calls are resolved
on the first attempt. On a busy afternoon, 80 calls come in. Assume $X \sim B(80, 0.22)$.

\vspace{8pt}
(a) Calculate $\mu_X$ and $\sigma_X$. Show formula setup.

\vspace{4pt}
$\mu_X = $ \blank{6cm} $= $ \blank{2cm}

$\sigma_X = $ \blank{6cm} $\approx$ \blank{2cm}

\vspace{8pt}
(b) Interpret $\mu_X$ in context.

\vspace{4pt}
\writeline\\[1pt]\writeline

\vspace{8pt}
(c) Find $P(X \le 15)$. Write the calculator notation and the answer.

\vspace{4pt}
$P(X \le 15) = \texttt{binomcdf(}\blank{1cm}\texttt{,}\;\blank{1cm}\texttt{,}\;15\texttt{)} \approx $ \blank{2cm}

\vspace{8pt}
(d) Find $P(X \ge 20)$. Write the inequality rewrite, then the calculator notation.

\vspace{4pt}
$P(X \ge 20) = 1 - P(X \le \blank{1.5cm}) = 1 - \texttt{binomcdf(}\blank{1cm}\texttt{,}\;\blank{1cm}\texttt{,}\;\blank{1cm}\texttt{)} \approx $ \blank{2cm}

\vspace{8pt}
(e) Is 20 or more resolved calls on the first attempt \emph{unusually high}?
Use $\mu_X$ and $\sigma_X$ to support your answer.

\vspace{4pt}
\writeline\\[1pt]\writeline
\vspace{4pt}
\end{tcolorbox}

\vspace{0.1in}

% ── TIER E ────────────────────────────────────────────────────────────────────
\begin{tcolorbox}[colback=greenbg, colframe=greenacc, boxrule=1pt, arc=2mm,
    title={\bfseries Tier E \quad Extension: Comparing Two Binomials},
    fonttitle=\small]
\small

\textbf{Context:} Two quality-control processes are being evaluated.
\begin{itemize}[itemsep=3pt]
  \item Process A: $n = 200$ items inspected, defect rate $p = 0.05$.
  \item Process B: $n = 100$ items inspected, defect rate $p = 0.10$.
\end{itemize}

\vspace{6pt}
(a) Compute $\mu_X$ and $\sigma_X$ for both processes.

\vspace{4pt}
\renewcommand{\arraystretch}{1.8}
\begin{tabularx}{\linewidth}{@{} >{\bfseries}p{2cm} X X @{}}
 & Process A & Process B \\
\hline
$\mu_X$ & \blank{4cm} & \blank{4cm} \\
$\sigma_X$ & \blank{4cm} & \blank{4cm} \\
\end{tabularx}

\vspace{8pt}
(b) Both processes have the same mean. Explain why in terms of the formulas.

\vspace{4pt}
\writeline\\[1pt]\writeline

\vspace{8pt}
(c) Which process has more variability? What does this mean practically for the
quality-control engineer?

\vspace{4pt}
\writeline\\[1pt]\writeline

\vspace{8pt}
(d) Challenge: For a fixed $n$, what value of $p$ maximizes $\sigma_X$?
Show a brief algebraic or numerical argument.

\vspace{4pt}
\writeline\\[1pt]\writeline\\[1pt]\writeline
\vspace{4pt}
\end{tcolorbox}

\end{document}
